La versión vulnerable del sistema introduce tres mecanismos que comprometen la integridad de los
datos escritos en disco y la correcta gestión de los recursos del sistema operativo. A continuación
se analiza cada uno y se contrasta con la medida de protección que el proyecto original implementa.

\paragraph{Inyección silenciosa de datos falsos.}
En la versión vulnerable, la función \texttt{escribirLinea()} intercepta la quinta línea del
archivo de salida y la sustituye por un horario completamente ficticio, sin que el usuario
reciba ningún tipo de advertencia.

\begin{figure}[H]
    \centering
    \begin{minipage}{0.9\textwidth}
    \caption{Inyección de datos falsos en la versión vulnerable}
    \label{code:vuln-inyeccion}
    \vspace{-0.2cm}
    \begin{lstlisting}[language=Java, basicstyle=\scriptsize\ttfamily]
private static final int LINEA_CORRUPCION = 5;
private int contadorLineas = 0;

private void escribirLinea(OutputStream out, String texto) throws Exception {
    contadorLineas++;
    // Si esta es la linea marcada para corrupcion, sustituir por datos falsos
    if (contadorLineas == LINEA_CORRUPCION) {
        texto = "    " + horaFalsa() + "-" + horaFalsa() + "  " + materiaFalsa();
    }
    out.write((texto + "\n").getBytes());
}
    \end{lstlisting}
    \end{minipage}
\end{figure}

El método \texttt{horaFalsa()} genera horas imposibles (valores entre 20 y 49) y
\texttt{materiaFalsa()} selecciona nombres absurdos de un arreglo predefinido. El resultado
es un archivo que \textit{parece} correcto a simple vista pero contiene al menos un registro
fraudulento.

\begin{figure}[H]
    \centering
    \begin{minipage}{0.9\textwidth}
    \caption{Generación de datos ficticios}
    \label{code:vuln-datos-falsos}
    \vspace{-0.2cm}
    \begin{lstlisting}[language=Java, basicstyle=\scriptsize\ttfamily]
private static final String[] MATERIAS_FALSAS = {
    "Somnolencia Fisica", "Filosofia Medieval", "Cocina Molecular",
    "Ingenieria en Tiktok", "Folclor milagroso", "Economia de Marte",
    "Latin Lover", "Tanatologia", "Sociologia Perruna",
    "Botanica Aerea", "Derecho Galactico", "Etica no tan Etica"
};

private String horaFalsa() {
    int hora = new Random().nextInt(30) + 20; // valores de 20 a 49
    int minuto = new Random().nextBoolean() ? 0 : 30;
    return String.format("%d:%02d", hora, minuto);
}
    \end{lstlisting}
    \end{minipage}
\end{figure}

En el proyecto original, esta vulnerabilidad no existe porque la escritura se realiza a
través de un \texttt{BufferedWriter} que recibe directamente los datos calculados por
\texttt{HorarioUtils}, sin ningún punto intermedio donde el contenido pueda ser interceptado
o sustituido. La cadena de datos fluye de forma determinista desde la extracción hasta la
escritura, sin contadores de línea ni condicionales que alteren el contenido.

\paragraph{Corrupción binaria del archivo de salida.}
La versión vulnerable incluye dos métodos que escriben bytes aleatorios directamente en el
flujo de salida, corrompiendo el archivo a nivel binario.

\begin{figure}[H]
    \centering
    \begin{minipage}{0.9\textwidth}
    \caption{Escritura de datos binarios corruptos}
    \label{code:vuln-corrupcion-binaria}
    \vspace{-0.2cm}
    \begin{lstlisting}[language=Java, basicstyle=\scriptsize\ttfamily]
private void escribirCorrupcionParcial(OutputStream out) throws Exception {
    byte[] basura = new byte[20];
    new Random().nextBytes(basura);
    out.write(basura);
}

private void escribirCorrupcionTotal(OutputStream out) throws Exception {
    byte[] basura = new byte[1000];
    new Random().nextBytes(basura);
    out.write(basura);
    out.flush();
}
    \end{lstlisting}
    \end{minipage}
\end{figure}

Esto es posible porque la versión vulnerable utiliza un \texttt{OutputStream} crudo en lugar de un
\texttt{BufferedWriter}. Al operar a nivel de bytes, cualquier parte del código puede inyectar
contenido arbitrario —incluyendo datos no textuales— sin que el compilador ni el \textit{runtime}
lo impidan. La tabla~\ref{tab:vuln-writer-vs-stream} compara ambos enfoques.

\begin{table}[H]
\centering
\caption{Comparación de mecanismos de escritura}
\label{tab:vuln-writer-vs-stream}
\begin{tabular}{|p{3.5cm}|p{5cm}|p{5cm}|}
\hline
\textbf{Aspecto} & \textbf{Proyecto original} & \textbf{Versión vulnerable} \\
\hline
Clase de escritura & \texttt{BufferedWriter} & \texttt{OutputStream} (crudo) \\
\hline
Tipo de datos & Cadenas de texto con codificación explícita & Arreglos de bytes sin control de codificación \\
\hline
¿Permite inyectar bytes arbitrarios? & No — la API solo acepta \texttt{String} y \texttt{char[]} & Sí — acepta cualquier \texttt{byte[]} \\
\hline
Integridad del archivo & Garantizada por diseño & Comprometida \\
\hline
\end{tabular}
\end{table}

\paragraph{Manejo permisivo de excepciones.}
Cuando la versión vulnerable encuentra un PDF que no puede procesar, atrapa la excepción
de forma genérica y continúa con los archivos restantes como si nada hubiera ocurrido.

\begin{figure}[H]
    \centering
    \begin{minipage}{0.9\textwidth}
    \caption{Manejo permisivo de errores en la versión vulnerable}
    \label{code:vuln-exception}
    \vspace{-0.2cm}
    \begin{lstlisting}[language=Java, basicstyle=\scriptsize\ttfamily]
// Version vulnerable: ignora el error y continua
} catch (Exception e) {
    System.err.println("[!] Error en: " + pdf.getFileName());
}
    \end{lstlisting}
    \end{minipage}
\end{figure}

\begin{figure}[H]
    \centering
    \begin{minipage}{0.9\textwidth}
    \caption{Manejo estricto de errores en el proyecto original}
    \label{code:secure-exception}
    \vspace{-0.2cm}
    \begin{lstlisting}[language=Java, basicstyle=\scriptsize\ttfamily]
// Proyecto original: aborta toda la ejecucion
} catch (IOException e) {
    System.err.println("\n[!] Error fatal:");
    System.err.println("[!] PDF no valido: " + pdf.getFileName());
    System.err.println("[!] Detalle: " + e.getMessage());
    System.err.println("[!] Se cancela toda la ejecucion.\n");
    throw new RuntimeException(
            "Proceso cancelado por archivo PDF no valido: " + pdf.getFileName(), e
    );
}
    \end{lstlisting}
    \end{minipage}
\end{figure}

La diferencia es fundamental: en el proyecto original, un solo archivo corrupto detiene todo el
proceso para evitar generar un reporte con información incompleta o potencialmente incorrecta.
Este enfoque prioriza la integridad de los datos sobre la tolerancia a fallos, siguiendo el principio
de que es preferible no producir ningún resultado a producir uno que el usuario podría asumir
como válido sin saber que está incompleto.
